% Options for packages loaded elsewhere
\PassOptionsToPackage{unicode}{hyperref}
\PassOptionsToPackage{hyphens}{url}
\documentclass[
]{book}
\usepackage{xcolor}
\usepackage{amsmath,amssymb}
\setcounter{secnumdepth}{5}
\usepackage{iftex}
\ifPDFTeX
  \usepackage[T1]{fontenc}
  \usepackage[utf8]{inputenc}
  \usepackage{textcomp} % provide euro and other symbols
\else % if luatex or xetex
  \usepackage{unicode-math} % this also loads fontspec
  \defaultfontfeatures{Scale=MatchLowercase}
  \defaultfontfeatures[\rmfamily]{Ligatures=TeX,Scale=1}
\fi
\usepackage{lmodern}
\ifPDFTeX\else
  % xetex/luatex font selection
\fi
% Use upquote if available, for straight quotes in verbatim environments
\IfFileExists{upquote.sty}{\usepackage{upquote}}{}
\IfFileExists{microtype.sty}{% use microtype if available
  \usepackage[]{microtype}
  \UseMicrotypeSet[protrusion]{basicmath} % disable protrusion for tt fonts
}{}
\makeatletter
\@ifundefined{KOMAClassName}{% if non-KOMA class
  \IfFileExists{parskip.sty}{%
    \usepackage{parskip}
  }{% else
    \setlength{\parindent}{0pt}
    \setlength{\parskip}{6pt plus 2pt minus 1pt}}
}{% if KOMA class
  \KOMAoptions{parskip=half}}
\makeatother
\usepackage{color}
\usepackage{fancyvrb}
\newcommand{\VerbBar}{|}
\newcommand{\VERB}{\Verb[commandchars=\\\{\}]}
\DefineVerbatimEnvironment{Highlighting}{Verbatim}{commandchars=\\\{\}}
% Add ',fontsize=\small' for more characters per line
\usepackage{framed}
\definecolor{shadecolor}{RGB}{248,248,248}
\newenvironment{Shaded}{\begin{snugshade}}{\end{snugshade}}
\newcommand{\AlertTok}[1]{\textcolor[rgb]{0.94,0.16,0.16}{#1}}
\newcommand{\AnnotationTok}[1]{\textcolor[rgb]{0.56,0.35,0.01}{\textbf{\textit{#1}}}}
\newcommand{\AttributeTok}[1]{\textcolor[rgb]{0.13,0.29,0.53}{#1}}
\newcommand{\BaseNTok}[1]{\textcolor[rgb]{0.00,0.00,0.81}{#1}}
\newcommand{\BuiltInTok}[1]{#1}
\newcommand{\CharTok}[1]{\textcolor[rgb]{0.31,0.60,0.02}{#1}}
\newcommand{\CommentTok}[1]{\textcolor[rgb]{0.56,0.35,0.01}{\textit{#1}}}
\newcommand{\CommentVarTok}[1]{\textcolor[rgb]{0.56,0.35,0.01}{\textbf{\textit{#1}}}}
\newcommand{\ConstantTok}[1]{\textcolor[rgb]{0.56,0.35,0.01}{#1}}
\newcommand{\ControlFlowTok}[1]{\textcolor[rgb]{0.13,0.29,0.53}{\textbf{#1}}}
\newcommand{\DataTypeTok}[1]{\textcolor[rgb]{0.13,0.29,0.53}{#1}}
\newcommand{\DecValTok}[1]{\textcolor[rgb]{0.00,0.00,0.81}{#1}}
\newcommand{\DocumentationTok}[1]{\textcolor[rgb]{0.56,0.35,0.01}{\textbf{\textit{#1}}}}
\newcommand{\ErrorTok}[1]{\textcolor[rgb]{0.64,0.00,0.00}{\textbf{#1}}}
\newcommand{\ExtensionTok}[1]{#1}
\newcommand{\FloatTok}[1]{\textcolor[rgb]{0.00,0.00,0.81}{#1}}
\newcommand{\FunctionTok}[1]{\textcolor[rgb]{0.13,0.29,0.53}{\textbf{#1}}}
\newcommand{\ImportTok}[1]{#1}
\newcommand{\InformationTok}[1]{\textcolor[rgb]{0.56,0.35,0.01}{\textbf{\textit{#1}}}}
\newcommand{\KeywordTok}[1]{\textcolor[rgb]{0.13,0.29,0.53}{\textbf{#1}}}
\newcommand{\NormalTok}[1]{#1}
\newcommand{\OperatorTok}[1]{\textcolor[rgb]{0.81,0.36,0.00}{\textbf{#1}}}
\newcommand{\OtherTok}[1]{\textcolor[rgb]{0.56,0.35,0.01}{#1}}
\newcommand{\PreprocessorTok}[1]{\textcolor[rgb]{0.56,0.35,0.01}{\textit{#1}}}
\newcommand{\RegionMarkerTok}[1]{#1}
\newcommand{\SpecialCharTok}[1]{\textcolor[rgb]{0.81,0.36,0.00}{\textbf{#1}}}
\newcommand{\SpecialStringTok}[1]{\textcolor[rgb]{0.31,0.60,0.02}{#1}}
\newcommand{\StringTok}[1]{\textcolor[rgb]{0.31,0.60,0.02}{#1}}
\newcommand{\VariableTok}[1]{\textcolor[rgb]{0.00,0.00,0.00}{#1}}
\newcommand{\VerbatimStringTok}[1]{\textcolor[rgb]{0.31,0.60,0.02}{#1}}
\newcommand{\WarningTok}[1]{\textcolor[rgb]{0.56,0.35,0.01}{\textbf{\textit{#1}}}}
\usepackage{longtable,booktabs,array}
\usepackage{calc} % for calculating minipage widths
% Correct order of tables after \paragraph or \subparagraph
\usepackage{etoolbox}
\makeatletter
\patchcmd\longtable{\par}{\if@noskipsec\mbox{}\fi\par}{}{}
\makeatother
% Allow footnotes in longtable head/foot
\IfFileExists{footnotehyper.sty}{\usepackage{footnotehyper}}{\usepackage{footnote}}
\makesavenoteenv{longtable}
\usepackage{graphicx}
\makeatletter
\newsavebox\pandoc@box
\newcommand*\pandocbounded[1]{% scales image to fit in text height/width
  \sbox\pandoc@box{#1}%
  \Gscale@div\@tempa{\textheight}{\dimexpr\ht\pandoc@box+\dp\pandoc@box\relax}%
  \Gscale@div\@tempb{\linewidth}{\wd\pandoc@box}%
  \ifdim\@tempb\p@<\@tempa\p@\let\@tempa\@tempb\fi% select the smaller of both
  \ifdim\@tempa\p@<\p@\scalebox{\@tempa}{\usebox\pandoc@box}%
  \else\usebox{\pandoc@box}%
  \fi%
}
% Set default figure placement to htbp
\def\fps@figure{htbp}
\makeatother
\setlength{\emergencystretch}{3em} % prevent overfull lines
\providecommand{\tightlist}{%
  \setlength{\itemsep}{0pt}\setlength{\parskip}{0pt}}
\usepackage[]{natbib}
\bibliographystyle{apalike}
\usepackage{booktabs}
\usepackage{amsthm}
\makeatletter
\def\thm@space@setup{%
  \thm@preskip=8pt plus 2pt minus 4pt
  \thm@postskip=\thm@preskip
}
\makeatother
\usepackage{bookmark}
\IfFileExists{xurl.sty}{\usepackage{xurl}}{} % add URL line breaks if available
\urlstyle{same}
\hypersetup{
  pdftitle={Introducción a R para la asignatura Probabilidad \& Estadística},
  pdfauthor={Daniel Pérez Troncoso},
  hidelinks,
  pdfcreator={LaTeX via pandoc}}

\title{Introducción a R para la asignatura Probabilidad \& Estadística}
\author{Daniel Pérez Troncoso}
\date{2026-03-09}

\begin{document}
\maketitle

{
\setcounter{tocdepth}{1}
\tableofcontents
}
\chapter{Prerrequisitos}\label{prerrequisitos}

\begin{enumerate}
\def\labelenumi{\arabic{enumi}.}
\tightlist
\item
  Instala R desde \href{https://cran.r-project.org/}{CRAN}
\item
  Instala RStudio desde \href{https://posit.co/download/rstudio-desktop/}{POSIT}
\item
  Abre RStudio por primera vez y crea un archivo .r: File / New File / R Script
\item
  Decide donde guardar el .r: File / Save as
\item
  Ese .r es una `receta' del código que se ejecutará en la consola de R.
\item
  A partir de ahora podrás escribir código en tu .r. Para ejecutarlo selecciona la línea que quieres ejecutar y pulsa Ctrl + Enter o ejecuta todo el código de una vez en \texttt{Run} arriba a la derecha de la ventana donde leerás el .r
\end{enumerate}

Algo fundamental es conocer el \emph{directorio de trabajo}. R se ejecuta sobre una localización de tu ordenador en concreto. Si quieres importar archivos a R (como, por ejemplo, una base de datos en excel), lo mejor es que lo pongas en el mismo directorio en el que se está ejecutando R. Con el comando \texttt{getwd()} \emph{(``get the working directory'')}, podrás conocer el directorio.

\begin{Shaded}
\begin{Highlighting}[]
\FunctionTok{getwd}\NormalTok{()}
\end{Highlighting}
\end{Shaded}

La recomendación es siempre establecer el directorio de trabajo en el mismo lugar en el que tengas guardado el archivo \emph{.r}. Esto ocurre automáticamente siempre que abres R haciendo doble clic sobre un archivo .r. Si tienes algún problema puedes ir al menú superior y hacer \texttt{Session\ /\ Set\ Working\ Directory\ /\ To\ Source\ File\ Location}.

\begin{quote}
Consejo: cada vez que escribas una nueva línea de código en tu archivo .r recuerda hacer \texttt{Ctrl\ +\ R} para guardar los cambios.
\end{quote}

\chapter{Cargando paquetes y datos}\label{cpyd}

Los paquetes contienen funciones que nos permiten cargar archivos excel, representar datos, o hacer ciertas operaciones. R es un programa de código libre al que contribuyen sus usuarios, por lo que hay miles de paquetes a nuestra disposición.

En esta clase solo haremos uso de dos paquetes: uno para cargar un archivo excel, ``readxl'' y otro para representar gráficos ``ggplot2''. Para ello, es necesario que primero instalemos y después activemos los paquetes paquetes.

Si los paquetes no están instalados, lo que es habitual la primera vez que se ejecuta R, deberemos instalarlos con:

\begin{Shaded}
\begin{Highlighting}[]
\CommentTok{\# Aquellas líneas con \# delante son comentarios. Ejecuta las siguientes líneas sin \#.}
\CommentTok{\# install.packages("readxl")}
\CommentTok{\# install.packages("ggplot2")}
\end{Highlighting}
\end{Shaded}

Una vez que los paquetes han sido ejecutados, podemos activarlos para avisar a R de que necesitamos cargar las funciones que emplean estos paquetes:

\begin{Shaded}
\begin{Highlighting}[]
\FunctionTok{library}\NormalTok{(readxl)}
\FunctionTok{library}\NormalTok{(ggplot2)}
\end{Highlighting}
\end{Shaded}

\section{Cargando datos}\label{cargando-datos}

Lo más sencillo es trabajar con directorios relativos. Esto implica colocar el archivo .r en el mismo directorio que el archivo de datos que quiere importarse. Por ejemplo, la estructura de datos que debes montar es:

\begin{verbatim}
📁 Mi análisis con R
|
|__ 📃 analisis.r
|
|__ 📃 datos.xlsx
\end{verbatim}

Una vez la estructura de datos sea así, podemos importar el excel como:

\begin{Shaded}
\begin{Highlighting}[]
\NormalTok{datos\_mkt }\OtherTok{\textless{}{-}} \FunctionTok{read\_excel}\NormalTok{(}\AttributeTok{path =} \StringTok{"datos\_mkt.xlsx"}\NormalTok{, }\AttributeTok{sheet =} \DecValTok{3}\NormalTok{)}
\NormalTok{datos\_econ }\OtherTok{\textless{}{-}} \FunctionTok{read\_excel}\NormalTok{(}\AttributeTok{path =} \StringTok{"datos\_econ.xlsx"}\NormalTok{, }\AttributeTok{sheet =} \DecValTok{2}\NormalTok{)}
\CommentTok{\# Eso es lo mismo que: }
\CommentTok{\# datos\_econ \textless{}{-} read\_excel("datos\_econ.xlsx", 2)}
\CommentTok{\# No hace falta que pongamos el nombre de los argumentos siempre que respetemos el orden. }
\end{Highlighting}
\end{Shaded}

\section{Echando un vistazo a nuestros datos}\label{echando-un-vistazo-a-nuestros-datos}

Al ejecutar la línea anterior, el objeto ``datos'' debe de aparecer en nuestro entorno. Vamos a echar un vistazo y ver qué hay dentro; la función \texttt{head()} nos muestra las primeras filas. También podemos pulsar sobre ``datos'' en el entorno y se abrirá una pestaña mostrándonos su contenido.

\begin{Shaded}
\begin{Highlighting}[]
\FunctionTok{head}\NormalTok{(datos\_econ)}
\end{Highlighting}
\end{Shaded}

\begin{verbatim}
## # A tibble: 6 x 2
##   Country  `GDP (current, M€)`
##   <chr>                  <dbl>
## 1 Belgium              620272.
## 2 Bulgaria             104767.
## 3 Czechia              320786.
## 4 Denmark              392401.
## 5 Germany             4328970 
## 6 Estonia               39848.
\end{verbatim}

\begin{Shaded}
\begin{Highlighting}[]
\FunctionTok{head}\NormalTok{(datos\_mkt)}
\end{Highlighting}
\end{Shaded}

\begin{verbatim}
## # A tibble: 6 x 17
##     age job        marital education default balance housing loan  contact   day
##   <dbl> <chr>      <chr>   <chr>     <chr>     <dbl> <chr>   <chr> <chr>   <dbl>
## 1    30 unemployed married primary   no         1787 no      no    cellul~    19
## 2    33 services   married secondary no         4789 yes     yes   cellul~    11
## 3    35 management single  tertiary  no         1350 yes     no    cellul~    16
## 4    30 management married tertiary  no         1476 yes     yes   unknown     3
## 5    59 blue-coll~ married secondary no            0 yes     no    unknown     5
## 6    35 management single  tertiary  no          747 no      no    cellul~    23
## # i 7 more variables: month <chr>, duration <dbl>, campaign <dbl>, pdays <dbl>,
## #   previous <dbl>, poutcome <chr>, y <chr>
\end{verbatim}

\chapter{Manipulando datos}\label{manipulando_datos}

\begin{Shaded}
\begin{Highlighting}[]
\FunctionTok{head}\NormalTok{(datos\_econ)}
\end{Highlighting}
\end{Shaded}

\begin{verbatim}
## # A tibble: 6 x 2
##   Country  `GDP (current, M€)`
##   <chr>                  <dbl>
## 1 Belgium              620272.
## 2 Bulgaria             104767.
## 3 Czechia              320786.
## 4 Denmark              392401.
## 5 Germany             4328970 
## 6 Estonia               39848.
\end{verbatim}

\begin{Shaded}
\begin{Highlighting}[]
\FunctionTok{head}\NormalTok{(datos\_mkt)}
\end{Highlighting}
\end{Shaded}

\begin{verbatim}
## # A tibble: 6 x 17
##     age job        marital education default balance housing loan  contact   day
##   <dbl> <chr>      <chr>   <chr>     <chr>     <dbl> <chr>   <chr> <chr>   <dbl>
## 1    30 unemployed married primary   no         1787 no      no    cellul~    19
## 2    33 services   married secondary no         4789 yes     yes   cellul~    11
## 3    35 management single  tertiary  no         1350 yes     no    cellul~    16
## 4    30 management married tertiary  no         1476 yes     yes   unknown     3
## 5    59 blue-coll~ married secondary no            0 yes     no    unknown     5
## 6    35 management single  tertiary  no          747 no      no    cellul~    23
## # i 7 more variables: month <chr>, duration <dbl>, campaign <dbl>, pdays <dbl>,
## #   previous <dbl>, poutcome <chr>, y <chr>
\end{verbatim}

Vemos los datos anteriores contienen diversas observaciones (filas) y variables (columnas). Pero los nombres de las columnas están en inglés y además son demasiado complejos para manejarlo en un programa estadístico. Vamos a simplificar los nombres:

\begin{Shaded}
\begin{Highlighting}[]
\FunctionTok{names}\NormalTok{(datos\_econ) }\OtherTok{\textless{}{-}} \FunctionTok{c}\NormalTok{(}\StringTok{"pais"}\NormalTok{, }\StringTok{"pib"}\NormalTok{)}
\FunctionTok{names}\NormalTok{(datos\_mkt) }\OtherTok{\textless{}{-}} \FunctionTok{c}\NormalTok{(}\StringTok{"edad"}\NormalTok{, }\StringTok{"trab"}\NormalTok{, }\StringTok{"civil"}\NormalTok{ ,}\StringTok{"educ"}\NormalTok{, }\StringTok{"def"}\NormalTok{, }\StringTok{"balance"}\NormalTok{, }\StringTok{"casa"}\NormalTok{, }\StringTok{"prestam"}\NormalTok{, }\StringTok{"contrato"}\NormalTok{, }\StringTok{"dia"}\NormalTok{, }\StringTok{"mes"}\NormalTok{, }\StringTok{"durac"}\NormalTok{, }\StringTok{"camp"}\NormalTok{, }\StringTok{"pdias"}\NormalTok{, }\StringTok{"previo"}\NormalTok{, }\StringTok{"pres"}\NormalTok{, }\StringTok{"y"}\NormalTok{)}
\CommentTok{\# La función names() extrae los nombres de \textquotesingle{}datos\textquotesingle{}. Usando \textless{}{-} estoy introduciendo los nombres los encabezados de \textquotesingle{}datos\textquotesingle{}.}
\FunctionTok{head}\NormalTok{(datos\_mkt)}
\end{Highlighting}
\end{Shaded}

\begin{verbatim}
## # A tibble: 6 x 17
##    edad trab  civil educ  def   balance casa  prestam contrato   dia mes   durac
##   <dbl> <chr> <chr> <chr> <chr>   <dbl> <chr> <chr>   <chr>    <dbl> <chr> <dbl>
## 1    30 unem~ marr~ prim~ no       1787 no    no      cellular    19 oct      79
## 2    33 serv~ marr~ seco~ no       4789 yes   yes     cellular    11 may     220
## 3    35 mana~ sing~ tert~ no       1350 yes   no      cellular    16 apr     185
## 4    30 mana~ marr~ tert~ no       1476 yes   yes     unknown      3 jun     199
## 5    59 blue~ marr~ seco~ no          0 yes   no      unknown      5 may     226
## 6    35 mana~ sing~ tert~ no        747 no    no      cellular    23 feb     141
## # i 5 more variables: camp <dbl>, pdias <dbl>, previo <dbl>, pres <chr>,
## #   y <chr>
\end{verbatim}

\begin{Shaded}
\begin{Highlighting}[]
\FunctionTok{head}\NormalTok{(datos\_econ)}
\end{Highlighting}
\end{Shaded}

\begin{verbatim}
## # A tibble: 6 x 2
##   pais          pib
##   <chr>       <dbl>
## 1 Belgium   620272.
## 2 Bulgaria  104767.
## 3 Czechia   320786.
## 4 Denmark   392401.
## 5 Germany  4328970 
## 6 Estonia    39848.
\end{verbatim}

\begin{Shaded}
\begin{Highlighting}[]
\CommentTok{\# También podemos escribir \textquotesingle{}datos\_mkt\textquotesingle{} o \textquotesingle{}datos\_econ\textquotesingle{} directamente en la consola para que me muestre el contenido del objeto.}
\end{Highlighting}
\end{Shaded}

¿Pero qué es la función \texttt{c()} que estoy usando? ¡Fácil! \texttt{c()} simplemente nos ayuda a crear otro objeto: un vector de datos, algo así como una fila individual de datos. Voy a crear uno que se llame \texttt{info} que contenga información sobre mi:

\begin{Shaded}
\begin{Highlighting}[]
\CommentTok{\# Voy a crear un vector de dos elementos donde el primero sea mi nombre y el segundo mi profesión. }
\NormalTok{info }\OtherTok{\textless{}{-}} \FunctionTok{c}\NormalTok{(}\StringTok{"Daniel"}\NormalTok{, }\StringTok{"Profesor"}\NormalTok{)}
\CommentTok{\# Si luego escribimos \textquotesingle{}info\textquotesingle{} en la consola se mostrará el contenido: }
\NormalTok{info}
\end{Highlighting}
\end{Shaded}

\begin{verbatim}
## [1] "Daniel"   "Profesor"
\end{verbatim}

De hecho, cada columna de \texttt{datos} puede extraerse como un vector, es muy sencillo. Simplemente hay que escribir \texttt{datos\$columna} donde columna deberá tener el nombre de la variable que queremos extraer:

\begin{Shaded}
\begin{Highlighting}[]
\NormalTok{vector\_pibs }\OtherTok{\textless{}{-}}\NormalTok{ datos\_econ}\SpecialCharTok{$}\NormalTok{pib}
\NormalTok{vector\_pibs}
\end{Highlighting}
\end{Shaded}

\begin{verbatim}
##  [1]  620271.8  104767.2  320786.5  392400.7 4328970.0   39847.7  562794.2
##  [8]  236736.0 1594330.0 2919899.9   85905.2 2199619.4   34770.2   40359.4
## [15]   78996.2   86180.3  206040.1   23136.7 1122459.0  494087.6  848490.9
## [22]  289784.3  353633.1   67418.1  130207.5  275993.0  558672.8   30696.5
## [29]  446866.1  896448.7   27483.1    7644.5   15672.0   25000.7   83257.9
## [36] 1253386.7
\end{verbatim}

\begin{Shaded}
\begin{Highlighting}[]
\NormalTok{vector\_edades }\OtherTok{\textless{}{-}}\NormalTok{ datos\_mkt}\SpecialCharTok{$}\NormalTok{edad}
\NormalTok{vector\_edades}
\end{Highlighting}
\end{Shaded}

\begin{verbatim}
##    [1] 30 33 35 30 59 35 36 39 41 43 39 43 36 20 31 40 56 37 25 31 38 42 44 44
##   [25] 26 41 55 67 56 53 68 31 59 32 49 42 78 32 33 23 38 36 52 32 32 34 55 26
##   [49] 32 61 45 37 38 34 53 48 57 33 36 54 41 63 48 48 56 51 31 29 41 32 37 42
##   [73] 31 41 56 54 41 40 50 40 27 60 28 52 37 51 41 35 38 34 39 34 27 32 57 41
##   [97] 41 30 36 31 36 38 49 57 34 32 37 33 56 27 21 58 32 33 25 25 30 22 39 46
##  [121] 31 45 32 36 51 34 24 45 32 77 41 37 43 48 27 33 29 34 52 60 32 32 56 31
##  [145] 42 35 56 48 46 37 41 36 45 55 39 35 55 55 27 34 34 25 41 43 28 31 78 45
##  [169] 48 32 42 37 40 32 42 43 44 37 37 37 28 32 34 40 34 34 34 37 36 56 49 33
##  [193] 30 36 23 57 27 38 42 75 34 48 35 39 29 45 42 29 38 56 36 25 38 57 31 42
##  [217] 45 30 49 49 53 37 49 41 54 32 33 33 36 35 52 40 44 45 30 48 33 57 57 51
##  [241] 34 50 44 27 50 36 42 32 42 43 31 32 55 58 46 55 34 28 42 29 54 32 39 53
##  [265] 58 40 29 33 38 56 40 24 46 35 35 34 35 38 50 51 59 32 37 42 40 28 47 27
##  [289] 41 34 50 47 53 53 47 49 29 45 38 32 70 65 56 35 56 56 49 38 35 38 49 55
##  [313] 43 43 25 30 46 50 35 26 45 37 60 38 47 33 51 35 35 48 37 29 34 32 30 47
##  [337] 33 33 63 51 41 56 60 41 37 36 57 55 44 47 31 35 39 50 29 30 46 27 48 26
##  [361] 36 42 50 49 45 23 38 32 60 30 58 60 38 51 42 30 40 37 53 27 39 39 59 37
##  [385] 41 45 48 48 32 35 33 28 64 35 33 41 34 37 38 39 35 28 27 48 50 45 41 39
##  [409] 44 54 49 34 75 42 53 49 31 45 39 46 42 37 26 37 36 29 36 26 32 40 53 29
##  [433] 34 43 56 55 31 27 36 44 37 30 42 30 62 24 31 66 47 39 35 35 27 32 44 42
##  [457] 34 25 40 43 55 50 45 35 32 59 34 39 40 36 40 55 38 61 41 57 48 77 35 54
##  [481] 37 56 40 40 54 55 55 52 55 33 52 33 53 46 44 48 50 34 38 42 32 58 31 19
##  [505] 41 31 52 42 41 30 38 32 52 48 36 45 34 36 34 61 54 50 43 47 45 57 27 49
##  [529] 43 40 32 58 34 46 29 52 45 26 40 41 47 27 27 48 49 35 47 54 39 28 54 34
##  [553] 63 42 43 37 55 43 41 47 34 28 36 50 35 40 26 36 59 28 46 50 53 81 35 58
##  [577] 44 38 49 33 29 28 31 33 52 43 32 36 41 43 45 39 41 31 37 28 47 37 60 43
##  [601] 58 52 50 32 42 31 28 36 51 46 50 32 26 41 39 31 28 31 30 35 33 30 32 40
##  [625] 31 36 33 57 45 59 30 33 52 83 58 33 41 31 35 27 57 56 53 39 32 46 56 46
##  [649] 29 60 33 27 50 45 58 59 58 64 47 40 36 28 38 53 46 54 32 55 55 56 40 44
##  [673] 32 33 41 37 58 52 49 44 28 31 37 30 28 46 47 35 80 30 55 22 33 53 49 51
##  [697] 41 34 36 35 29 40 32 48 59 35 33 55 41 42 46 35 28 25 56 30 45 40 40 39
##  [721] 48 31 32 34 26 31 37 29 57 47 38 48 50 33 28 30 52 36 36 46 57 28 46 55
##  [745] 33 53 41 60 35 43 57 46 31 53 35 32 48 50 40 42 59 27 32 35 28 37 57 42
##  [769] 55 52 26 52 26 33 24 37 28 30 33 55 33 42 37 59 27 23 51 50 43 37 34 41
##  [793] 26 49 29 30 38 58 34 53 50 30 40 57 40 41 53 48 32 34 46 45 51 31 47 32
##  [817] 41 37 27 27 50 31 40 43 39 31 48 32 47 35 42 47 59 38 50 27 55 39 31 37
##  [841] 28 42 40 59 40 43 38 54 41 55 33 46 36 48 42 46 49 36 34 51 58 48 41 49
##  [865] 43 32 54 33 35 56 57 31 33 29 38 37 33 30 31 28 26 31 50 45 48 41 60 41
##  [889] 43 36 40 28 27 49 37 31 31 31 35 51 33 30 58 34 45 28 31 40 30 34 40 59
##  [913] 43 41 33 33 36 32 25 25 38 34 67 32 45 33 38 32 54 51 42 41 39 38 39 46
##  [937] 50 33 37 31 34 34 36 40 40 41 52 40 48 33 44 55 51 34 58 39 27 55 35 37
##  [961] 52 41 52 34 53 33 47 35 33 55 58 44 25 37 41 34 36 44 34 56 43 28 54 37
##  [985] 59 58 56 32 42 52 46 37 39 31 59 38 58 34 40 20 34 31 39 34 24 41 40 45
## [1009] 36 53 23 42 48 34 36 43 36 49 55 53 42 45 31 32 36 56 46 41 31 32 38 49
## [1033] 46 29 36 38 53 43 41 42 55 39 30 41 44 58 35 27 34 25 36 38 30 47 41 36
## [1057] 63 23 32 57 48 49 50 33 53 30 38 44 39 45 36 38 36 30 28 26 54 38 32 50
## [1081] 43 48 57 64 26 33 50 47 39 52 31 50 47 30 32 36 46 45 37 58 59 28 48 55
## [1105] 37 41 35 34 43 59 50 37 35 50 36 41 28 42 60 28 48 51 32 35 37 38 77 32
## [1129] 43 36 29 32 41 26 34 57 36 42 35 52 36 46 60 57 48 59 39 44 27 59 36 44
## [1153] 50 41 47 32 38 55 38 31 38 29 55 35 40 34 32 33 33 37 58 36 42 34 35 60
## [1177] 50 28 27 37 31 34 38 32 37 33 71 27 38 34 54 30 35 33 36 59 72 64 53 35
## [1201] 58 40 41 48 44 39 61 36 28 58 60 66 36 49 40 23 29 44 54 47 33 30 27 49
## [1225] 27 51 32 30 33 27 75 37 36 37 35 32 32 35 31 28 27 22 35 69 37 45 47 52
## [1249] 33 60 33 44 38 54 32 33 55 51 48 37 33 30 43 49 54 31 54 36 30 27 30 30
## [1273] 35 33 43 41 41 46 54 32 30 40 30 27 41 54 32 31 42 48 34 55 49 46 47 24
## [1297] 58 32 25 49 59 30 46 42 37 44 49 41 48 31 37 51 80 32 48 49 30 27 34 49
## [1321] 48 40 38 46 41 28 51 33 26 34 27 42 45 31 32 39 26 40 52 40 48 31 27 47
## [1345] 28 40 34 29 36 83 55 43 23 44 58 29 30 29 41 35 31 27 47 53 34 34 36 40
## [1369] 37 65 53 28 33 56 45 39 28 36 41 46 41 59 24 34 40 42 50 25 51 50 37 21
## [1393] 52 55 26 38 32 31 34 33 35 44 48 29 50 37 44 36 49 31 49 35 57 42 36 75
## [1417] 38 40 38 34 29 28 79 51 32 36 38 49 42 43 34 59 35 33 34 64 34 77 57 30
## [1441] 56 49 73 39 57 32 32 45 32 31 49 29 55 30 45 39 36 33 36 31 36 38 33 40
## [1465] 31 39 33 53 38 59 55 38 48 53 44 57 29 31 35 31 35 40 39 43 44 29 40 33
## [1489] 46 27 26 41 34 39 27 31 44 55 51 50 30 46 49 56 47 30 36 40 37 42 45 34
## [1513] 27 30 32 46 29 53 32 49 35 38 39 34 42 35 60 57 31 29 41 46 33 32 54 37
## [1537] 32 43 50 34 31 43 49 22 35 56 53 54 32 41 40 33 33 34 36 34 55 28 29 31
## [1561] 52 54 43 55 28 45 45 51 42 43 34 36 37 50 42 34 48 47 40 37 38 29 46 26
## [1585] 30 42 42 31 37 26 54 70 29 50 60 32 32 34 48 25 32 32 47 32 49 53 41 47
## [1609] 35 33 49 46 29 58 40 34 60 43 50 33 37 55 45 47 58 36 55 31 38 42 32 39
## [1633] 41 59 42 33 53 43 39 27 39 27 33 29 33 41 38 30 58 37 36 45 23 53 44 32
## [1657] 32 38 47 34 52 45 50 52 37 51 32 33 28 66 31 35 47 36 32 34 40 33 39 38
## [1681] 32 65 34 34 42 31 47 44 38 28 46 51 42 37 58 43 36 34 31 49 54 29 51 40
## [1705] 35 25 32 30 38 46 43 32 29 56 57 38 51 35 37 48 30 60 56 45 38 20 59 34
## [1729] 36 53 51 58 48 56 51 35 40 59 44 49 35 41 55 51 27 57 51 38 49 40 47 48
## [1753] 47 56 29 38 32 42 30 26 35 35 35 57 29 32 45 51 31 50 38 37 35 37 28 38
## [1777] 43 43 56 32 58 47 31 38 50 32 37 39 40 41 45 31 54 39 56 35 34 31 38 41
## [1801] 44 49 30 31 34 70 45 31 57 40 28 44 49 46 34 58 52 27 48 33 46 51 31 35
## [1825] 37 60 40 50 34 32 36 32 40 37 57 46 31 39 35 32 44 41 59 39 54 27 51 33
## [1849] 38 52 54 33 50 39 34 41 48 38 34 32 48 28 47 32 32 38 86 47 42 44 25 34
## [1873] 34 53 42 30 26 40 34 34 40 36 34 32 39 36 39 33 34 42 46 32 27 35 53 29
## [1897] 35 49 58 40 19 30 34 60 28 46 49 59 44 33 54 26 54 48 30 29 51 53 43 34
## [1921] 40 43 33 66 44 33 31 71 22 46 32 45 32 49 27 46 37 30 44 45 27 57 66 40
## [1945] 48 33 42 47 27 78 50 46 41 36 49 46 77 30 43 44 40 38 25 49 33 40 44 32
## [1969] 52 35 35 32 42 59 44 50 27 55 54 31 56 51 61 34 56 29 34 35 56 45 31 37
## [1993] 43 40 44 38 58 24 53 36 28 45 32 46 28 51 51 30 39 28 31 31 38 31 74 40
## [2017] 43 34 54 51 33 41 33 48 42 35 25 38 45 59 47 33 59 41 42 51 45 30 33 43
## [2041] 37 73 43 52 44 56 21 43 54 59 52 31 53 24 57 48 57 51 50 34 43 31 29 39
## [2065] 36 38 56 41 47 38 80 37 44 33 53 69 46 30 31 38 34 48 46 42 41 33 39 27
## [2089] 26 38 48 54 51 47 53 35 58 52 57 52 35 48 46 39 38 34 47 35 46 27 31 31
## [2113] 46 31 33 42 33 40 51 31 62 24 48 41 43 52 54 32 38 41 31 35 26 52 58 57
## [2137] 34 59 27 42 47 33 31 37 29 38 35 31 62 32 31 51 29 43 36 34 29 57 43 33
## [2161] 60 51 40 51 38 35 47 45 40 29 52 26 31 55 35 56 28 27 30 51 33 32 43 46
## [2185] 32 40 39 51 46 43 36 48 26 39 47 23 54 38 33 35 52 25 29 35 49 41 32 29
## [2209] 28 59 34 55 34 24 42 34 34 51 58 44 33 34 34 42 46 31 39 34 47 42 48 30
## [2233] 38 24 34 46 32 46 26 42 50 41 34 32 58 45 28 35 39 49 50 37 45 30 48 36
## [2257] 41 55 70 41 47 34 33 58 51 45 55 30 33 27 59 34 47 44 35 45 42 38 34 35
## [2281] 29 58 44 45 43 44 30 38 50 21 38 42 45 34 48 34 32 35 39 36 35 57 28 43
## [2305] 57 34 36 47 36 33 57 57 34 51 54 52 43 58 53 44 29 27 35 40 34 38 32 41
## [2329] 25 38 26 37 30 36 33 33 35 24 29 54 27 34 39 38 37 26 36 35 30 25 34 52
## [2353] 39 34 34 30 33 30 51 36 48 37 58 56 34 33 26 40 35 45 37 35 43 28 32 26
## [2377] 35 35 39 42 32 30 26 30 32 47 33 57 35 53 24 54 31 41 44 35 30 43 51 60
## [2401] 43 31 56 42 52 40 36 40 46 32 31 41 58 36 52 42 32 62 31 41 34 44 42 37
## [2425] 47 32 46 39 28 26 52 42 33 73 35 40 51 37 39 28 49 36 45 29 35 35 30 46
## [2449] 25 36 42 57 42 29 36 31 59 50 47 46 26 31 30 47 23 44 42 32 28 57 38 22
## [2473] 54 32 32 59 45 53 50 34 35 30 62 40 26 34 50 31 35 55 34 58 50 31 47 38
## [2497] 46 42 34 37 38 59 32 55 35 56 27 42 42 41 53 67 53 33 31 31 30 48 30 43
## [2521] 38 46 59 44 35 55 39 36 33 48 33 67 48 57 31 30 56 26 44 44 31 47 30 37
## [2545] 56 43 42 60 34 34 41 38 23 33 36 49 34 40 72 43 44 35 46 55 50 23 34 40
## [2569] 61 53 57 28 25 38 27 45 31 39 42 46 36 30 32 52 31 60 35 38 46 45 59 34
## [2593] 33 46 48 31 32 35 36 52 45 34 25 26 57 48 48 32 58 57 31 33 53 36 32 45
## [2617] 35 41 54 49 28 24 30 32 53 42 48 33 31 56 58 37 59 50 39 30 47 44 38 50
## [2641] 27 73 46 56 36 33 40 52 51 33 54 44 50 45 36 42 37 28 56 59 32 57 49 56
## [2665] 48 27 31 28 36 58 46 32 39 35 31 29 48 60 74 30 52 24 29 33 34 45 31 45
## [2689] 27 30 45 57 34 40 39 46 55 25 34 37 64 45 57 21 58 41 41 58 33 32 44 47
## [2713] 41 54 47 32 59 29 48 32 46 31 53 34 48 32 31 27 60 34 52 48 45 54 52 47
## [2737] 54 34 39 46 26 37 47 35 57 34 47 27 40 39 36 42 23 34 34 49 57 59 57 34
## [2761] 61 35 42 57 42 53 63 47 54 45 28 58 48 35 37 37 37 25 58 52 19 59 41 34
## [2785] 35 41 28 53 35 33 37 32 32 54 44 36 56 30 36 43 40 32 34 37 28 35 53 60
## [2809] 40 57 29 44 41 47 43 37 34 30 57 26 34 46 38 37 28 57 39 49 38 27 30 38
## [2833] 36 38 59 37 29 43 33 54 55 32 48 33 46 44 47 38 80 40 43 32 39 33 31 36
## [2857] 35 25 33 23 31 36 44 31 55 51 36 27 40 49 59 32 42 35 29 29 37 49 51 61
## [2881] 43 42 24 46 23 26 43 35 38 56 31 32 28 26 49 33 80 40 42 46 32 37 32 53
## [2905] 39 31 34 44 35 39 55 34 39 58 37 56 32 41 29 44 24 43 26 39 49 52 33 30
## [2929] 32 27 58 46 28 39 41 38 50 56 34 33 61 36 34 36 45 27 45 40 36 46 37 36
## [2953] 37 22 51 49 35 50 31 47 59 32 71 31 41 35 36 46 46 47 35 37 56 46 44 48
## [2977] 28 38 34 56 55 44 38 53 41 34 34 35 38 42 56 37 47 29 41 43 38 31 32 40
## [3001] 45 43 27 33 37 60 51 44 48 57 29 50 36 49 49 35 39 34 57 37 25 43 37 36
## [3025] 49 38 39 40 34 39 53 44 32 50 29 55 73 52 49 63 34 32 49 40 53 53 43 35
## [3049] 35 25 31 48 40 46 30 47 23 29 38 56 52 55 43 34 43 55 34 34 43 39 39 34
## [3073] 55 67 34 47 33 45 36 59 33 52 53 41 45 60 32 28 31 53 58 38 46 53 56 52
## [3097] 38 37 37 37 66 44 61 34 33 30 27 37 38 35 50 39 35 36 28 39 39 30 32 28
## [3121] 47 53 57 47 56 38 52 46 57 42 57 51 53 37 47 48 39 42 36 34 61 32 35 54
## [3145] 31 45 32 52 36 31 53 43 47 46 34 55 23 75 58 35 27 24 48 31 57 53 42 50
## [3169] 59 38 39 28 55 33 36 38 28 32 35 58 30 33 36 32 34 42 41 63 38 54 39 46
## [3193] 37 76 28 33 46 30 60 56 24 58 79 32 33 77 60 54 30 50 50 37 41 49 41 40
## [3217] 38 34 49 72 34 31 33 31 55 25 34 53 39 41 32 29 56 19 39 32 43 32 53 52
## [3241] 43 48 27 27 51 39 33 50 59 58 54 60 38 39 40 55 36 31 38 52 45 32 58 35
## [3265] 36 55 45 49 44 37 35 36 44 32 56 32 35 44 59 31 60 33 45 45 57 36 30 30
## [3289] 59 58 49 48 28 59 40 32 58 28 27 38 53 60 33 52 64 40 52 28 49 28 29 87
## [3313] 37 36 35 33 60 25 69 41 48 35 30 59 32 28 60 36 32 34 58 53 31 33 46 29
## [3337] 37 58 34 34 49 53 42 40 30 31 55 36 31 26 31 49 54 47 61 49 51 53 35 25
## [3361] 79 53 21 33 40 34 52 27 57 51 31 49 32 28 33 50 42 45 34 56 33 34 37 37
## [3385] 35 38 32 52 49 46 42 36 32 49 46 40 36 36 34 46 32 32 32 43 25 55 24 40
## [3409] 38 48 36 57 45 33 38 26 30 55 34 60 36 65 39 52 55 26 40 52 29 41 30 28
## [3433] 26 58 29 48 35 26 49 34 32 31 48 40 42 48 38 53 59 28 41 50 31 30 59 33
## [3457] 35 30 71 30 24 34 34 52 33 38 41 42 38 48 34 46 33 40 36 32 33 55 58 25
## [3481] 47 40 59 40 30 25 58 39 40 38 66 47 37 36 50 76 58 42 58 28 41 30 32 27
## [3505] 48 36 54 36 52 34 43 32 33 51 39 44 49 25 39 54 32 36 54 31 45 31 45 33
## [3529] 46 43 50 44 42 26 47 48 60 27 30 28 58 37 36 39 52 55 30 36 39 47 42 42
## [3553] 58 68 43 44 49 57 61 36 31 51 31 56 30 38 35 31 39 50 39 26 55 43 34 51
## [3577] 31 32 53 44 38 47 36 30 46 33 43 38 27 38 41 29 39 42 47 30 34 45 43 59
## [3601] 50 39 36 44 36 27 29 70 44 55 36 28 50 49 71 31 45 42 45 47 60 35 31 33
## [3625] 50 51 34 56 29 31 44 55 35 36 54 32 36 60 32 57 34 27 52 45 58 39 48 36
## [3649] 31 41 30 57 29 29 30 44 33 56 26 52 32 36 26 37 34 26 32 30 27 42 33 37
## [3673] 45 36 39 33 47 49 46 31 43 30 43 42 30 33 41 29 38 29 80 38 36 38 44 31
## [3697] 31 57 37 50 60 51 48 42 36 43 31 49 35 44 34 46 49 32 35 51 57 59 40 49
## [3721] 51 35 29 32 37 49 45 34 53 42 43 52 34 34 27 43 41 72 35 41 36 51 37 42
## [3745] 26 23 50 37 35 37 79 39 47 44 36 33 42 35 31 58 46 41 58 33 31 55 38 55
## [3769] 47 51 46 61 47 34 34 41 36 26 35 36 26 58 50 61 28 49 74 53 46 47 42 56
## [3793] 42 49 43 39 66 30 38 31 32 48 34 28 40 36 52 46 25 52 40 26 26 35 40 31
## [3817] 32 33 53 46 28 31 44 35 50 32 51 51 31 41 57 34 47 47 36 29 39 53 47 51
## [3841] 48 32 36 50 39 31 49 33 26 27 51 26 45 54 49 47 42 50 37 54 40 50 57 33
## [3865] 49 35 28 58 31 36 49 30 29 50 26 27 32 62 34 25 34 44 38 57 36 41 36 27
## [3889] 37 47 31 39 35 40 38 31 36 53 59 29 30 52 55 56 32 24 51 36 36 42 49 48
## [3913] 36 48 30 27 38 36 45 44 42 37 48 57 44 48 56 37 47 54 51 47 25 39 40 38
## [3937] 34 53 32 35 24 58 49 45 55 42 28 48 51 38 27 50 44 40 29 34 44 50 34 54
## [3961] 34 34 52 60 39 58 41 48 58 40 48 52 50 25 54 34 48 37 66 32 45 43 41 40
## [3985] 44 43 30 32 57 31 49 31 41 44 54 37 34 43 44 44 53 36 58 26 34 55 55 34
## [4009] 41 38 28 41 30 39 41 36 37 59 26 43 42 46 31 43 48 41 32 45 52 37 38 36
## [4033] 58 37 30 47 40 42 71 54 53 33 30 39 32 34 39 75 33 40 37 69 52 31 28 29
## [4057] 46 28 41 42 34 28 43 32 33 70 36 53 30 45 31 28 45 51 53 40 53 31 30 53
## [4081] 40 26 47 51 45 35 36 57 42 32 30 36 43 69 31 35 26 46 30 33 34 48 31 46
## [4105] 40 39 33 50 84 40 55 40 46 43 34 35 57 25 32 52 31 34 42 47 50 30 33 56
## [4129] 34 33 32 34 37 30 29 56 31 59 36 51 54 37 32 42 29 52 40 30 55 49 33 32
## [4153] 21 42 56 39 30 55 43 53 57 39 56 29 38 59 42 46 31 35 50 42 36 25 42 27
## [4177] 32 57 48 34 33 26 50 60 47 31 52 36 43 30 32 31 29 48 32 30 48 51 34 35
## [4201] 49 36 37 37 49 44 40 57 39 38 51 27 38 30 36 37 53 38 48 26 33 47 45 32
## [4225] 69 65 37 45 36 62 37 39 38 31 58 44 31 33 36 31 45 30 34 34 39 41 32 32
## [4249] 36 35 32 33 49 35 31 42 34 32 29 32 50 58 59 25 57 30 26 50 39 45 37 47
## [4273] 50 40 35 53 35 54 35 31 46 52 34 46 54 37 39 50 44 56 40 36 33 33 42 30
## [4297] 33 26 26 37 59 49 33 46 55 49 51 27 44 40 36 46 30 35 31 37 49 27 32 70
## [4321] 22 58 33 83 50 53 37 55 49 41 40 48 27 30 37 37 40 73 50 53 47 33 50 51
## [4345] 42 34 44 32 60 60 29 39 47 30 31 34 51 38 47 46 33 30 40 38 63 33 59 29
## [4369] 54 57 57 28 30 28 45 29 42 31 59 41 57 43 30 26 52 32 39 31 83 34 40 39
## [4393] 50 30 43 49 41 35 57 29 48 38 41 39 31 51 35 45 28 57 36 49 34 65 26 52
## [4417] 32 33 43 26 45 30 37 45 34 35 41 42 48 45 47 37 32 40 35 31 36 44 52 28
## [4441] 45 38 39 28 31 46 22 49 44 31 53 55 35 35 37 50 37 37 29 31 40 47 55 33
## [4465] 53 38 61 43 42 36 31 59 29 33 44 43 42 32 43 39 23 43 46 31 38 53 37 49
## [4489] 31 45 58 35 32 28 26 48 30 31 31 45 38 34 42 60 42 32 42 33 42 51 36 46
## [4513] 40 49 38 32 33 57 57 28 44
\end{verbatim}

\chapter{Representación de datos en gráficos de barras}\label{representacion_datos_barras}

Para la representación de datos usaremos el paquete \texttt{ggplot2} dentro del cual encontraremos la función \texttt{ggplot}. La función \texttt{ggplot} es muy potente y muy usada por todos los usuarios de R, pero también tiene una curva de aprendizaje inclinada. Aquí vamos a dar unos trucos para que sea más sencillo emplearla.

\section{Consejo 1}\label{consejo-1}

\texttt{ggplot} funciona por capas. Primero establecemos los elementos básicos del gráfico y después le vamos añadiendo cosas.

\begin{Shaded}
\begin{Highlighting}[]
\CommentTok{\# Primero hay que escribir la función ggplot con los argumentos \textquotesingle{}data\textquotesingle{} y \textquotesingle{}mapping\textquotesingle{}.}
\FunctionTok{ggplot}\NormalTok{(}\AttributeTok{data =}\NormalTok{ datos\_econ, }\AttributeTok{mapping =} \FunctionTok{aes}\NormalTok{(}\AttributeTok{x =}\NormalTok{ pais, }\AttributeTok{y =}\NormalTok{ pib)) }\SpecialCharTok{+}
\CommentTok{\# La mayoría de usuarios prefiere escribir simplemente ggplot(datos, aes(x = pais, y = pib))}
  \FunctionTok{geom\_bar}\NormalTok{(}\AttributeTok{stat =} \StringTok{"identity"}\NormalTok{)}
\CommentTok{\# Geom bar indica que vamos a representar un diagrama de barras. Identity hace referencia al tipo de datos que tenemos.}

\CommentTok{\#En este ejemplo con los datos\_mkt geom\_bar con "count" hace el recuento de todos los estados civiles que hay y los pone en barras. }
\FunctionTok{ggplot}\NormalTok{(datos\_mkt, }\FunctionTok{aes}\NormalTok{(}\AttributeTok{x =}\NormalTok{ civil)) }\SpecialCharTok{+} \FunctionTok{geom\_bar}\NormalTok{(}\AttributeTok{stat =} \StringTok{"count"}\NormalTok{)}
\end{Highlighting}
\end{Shaded}

\begin{Shaded}
\begin{Highlighting}[]
\CommentTok{\# Vamos a guardar este gráfico en el objeto \textquotesingle{}grafico\textquotesingle{}}
\NormalTok{grafico\_econ }\OtherTok{\textless{}{-}} \FunctionTok{ggplot}\NormalTok{(datos\_econ, }\FunctionTok{aes}\NormalTok{(}\AttributeTok{x =}\NormalTok{ pais, }\AttributeTok{y =}\NormalTok{ pib)) }\SpecialCharTok{+} \FunctionTok{geom\_bar}\NormalTok{(}\AttributeTok{stat =} \StringTok{"identity"}\NormalTok{)}
\NormalTok{grafico\_econ}
\end{Highlighting}
\end{Shaded}

\pandocbounded{\includegraphics[keepaspectratio]{bookdown-demo_files/figure-latex/unnamed-chunk-10-1.pdf}}

\begin{Shaded}
\begin{Highlighting}[]
\NormalTok{grafico\_mkt }\OtherTok{\textless{}{-}} \FunctionTok{ggplot}\NormalTok{(datos\_mkt, }\FunctionTok{aes}\NormalTok{(}\AttributeTok{x =}\NormalTok{ civil)) }\SpecialCharTok{+} \FunctionTok{geom\_bar}\NormalTok{(}\AttributeTok{stat =} \StringTok{"count"}\NormalTok{)}
\NormalTok{grafico\_mkt}
\end{Highlighting}
\end{Shaded}

\pandocbounded{\includegraphics[keepaspectratio]{bookdown-demo_files/figure-latex/unnamed-chunk-10-2.pdf}}

Ahora tenemos un gráfico de barras, pero necesitamos varias cosas:
1. Poner cada barra en un color
2. Girar las etiquetas de los nombres para que puedan leerse bien en grafico\_econ
3. Personalizar las etiquetas en grafico\_mkt

Primero, hay que añadir el argumento \texttt{fill} a \texttt{aes} con la variable que queremos usar para el código de colores. Para evitar que se muestre una leyenda con todos los colores usamos el argumento \texttt{show.legend} igual a ``F'' (falso).

\begin{Shaded}
\begin{Highlighting}[]
\NormalTok{grafico\_econ }\OtherTok{\textless{}{-}} \FunctionTok{ggplot}\NormalTok{(datos\_econ, }\FunctionTok{aes}\NormalTok{(}\AttributeTok{x =}\NormalTok{ pais, }\AttributeTok{y =}\NormalTok{ pib, }\AttributeTok{fill =}\NormalTok{ pais)) }\SpecialCharTok{+} 
  \FunctionTok{geom\_bar}\NormalTok{(}\AttributeTok{stat =} \StringTok{"identity"}\NormalTok{, }\AttributeTok{show.legend =}\NormalTok{ F)}
\NormalTok{grafico\_econ}
\end{Highlighting}
\end{Shaded}

\pandocbounded{\includegraphics[keepaspectratio]{bookdown-demo_files/figure-latex/unnamed-chunk-11-1.pdf}}

\begin{Shaded}
\begin{Highlighting}[]
\NormalTok{grafico\_mkt }\OtherTok{\textless{}{-}} \FunctionTok{ggplot}\NormalTok{(datos\_mkt, }\FunctionTok{aes}\NormalTok{(}\AttributeTok{x =}\NormalTok{ civil, }\AttributeTok{fill =}\NormalTok{ civil)) }\SpecialCharTok{+} 
  \FunctionTok{geom\_bar}\NormalTok{(}\AttributeTok{stat =} \StringTok{"count"}\NormalTok{, }\AttributeTok{show.legend =}\NormalTok{ F)}
\NormalTok{grafico\_mkt}
\end{Highlighting}
\end{Shaded}

\pandocbounded{\includegraphics[keepaspectratio]{bookdown-demo_files/figure-latex/unnamed-chunk-11-2.pdf}}
Posteriormente, podemos sumar la función \texttt{theme} al gráfico para girar las etiquetas de los países:

\begin{Shaded}
\begin{Highlighting}[]
\NormalTok{grafico\_econ }\OtherTok{\textless{}{-}}\NormalTok{ grafico\_econ }\SpecialCharTok{+} \FunctionTok{theme}\NormalTok{(}\AttributeTok{axis.text.x =} \FunctionTok{element\_text}\NormalTok{(}\AttributeTok{angle =} \DecValTok{45}\NormalTok{, }\AttributeTok{hjust =} \DecValTok{1}\NormalTok{))}
\NormalTok{grafico\_econ}
\end{Highlighting}
\end{Shaded}

\pandocbounded{\includegraphics[keepaspectratio]{bookdown-demo_files/figure-latex/unnamed-chunk-12-1.pdf}}

En \texttt{datos\_mkt} debemos cambiar las etiquetas de \texttt{civil}si queremos que en el gráfico aparezcan las tres opciones bien etiquetadas (necesitamos modificar los valores de la variable en la base de datos):

\begin{Shaded}
\begin{Highlighting}[]
\NormalTok{datos\_mkt}\SpecialCharTok{$}\NormalTok{civil }\OtherTok{\textless{}{-}} \FunctionTok{ifelse}\NormalTok{(datos\_mkt}\SpecialCharTok{$}\NormalTok{civil }\SpecialCharTok{==} \StringTok{"married"}\NormalTok{, }\StringTok{"Casado/a"}\NormalTok{, datos\_mkt}\SpecialCharTok{$}\NormalTok{civil)}
\NormalTok{datos\_mkt}\SpecialCharTok{$}\NormalTok{civil }\OtherTok{\textless{}{-}} \FunctionTok{ifelse}\NormalTok{(datos\_mkt}\SpecialCharTok{$}\NormalTok{civil }\SpecialCharTok{==} \StringTok{"single"}\NormalTok{, }\StringTok{"Soltero/a"}\NormalTok{, datos\_mkt}\SpecialCharTok{$}\NormalTok{civil)}
\NormalTok{datos\_mkt}\SpecialCharTok{$}\NormalTok{civil }\OtherTok{\textless{}{-}} \FunctionTok{ifelse}\NormalTok{(datos\_mkt}\SpecialCharTok{$}\NormalTok{civil }\SpecialCharTok{==} \StringTok{"divorced"}\NormalTok{, }\StringTok{"Divorciado/a"}\NormalTok{, datos\_mkt}\SpecialCharTok{$}\NormalTok{civil)}
\NormalTok{grafico\_mkt }\OtherTok{\textless{}{-}} \FunctionTok{ggplot}\NormalTok{(datos\_mkt, }\FunctionTok{aes}\NormalTok{(}\AttributeTok{x =}\NormalTok{ civil, }\AttributeTok{fill =}\NormalTok{ civil)) }\SpecialCharTok{+} 
  \FunctionTok{geom\_bar}\NormalTok{(}\AttributeTok{stat =} \StringTok{"count"}\NormalTok{, }\AttributeTok{show.legend =}\NormalTok{ F)}
\NormalTok{grafico\_mkt}
\end{Highlighting}
\end{Shaded}

\pandocbounded{\includegraphics[keepaspectratio]{bookdown-demo_files/figure-latex/unnamed-chunk-13-1.pdf}}

Finalmente, con \texttt{labs()} podemos cambiar las etiquetas de los ejes:

\begin{Shaded}
\begin{Highlighting}[]
\NormalTok{grafico\_econ }\OtherTok{\textless{}{-}}\NormalTok{ grafico\_econ }\SpecialCharTok{+} \FunctionTok{labs}\NormalTok{(}\AttributeTok{x =} \StringTok{"Países"}\NormalTok{, }\AttributeTok{y =} \StringTok{"PIB (Millones de €)"}\NormalTok{)}
\NormalTok{grafico\_econ}
\end{Highlighting}
\end{Shaded}

\pandocbounded{\includegraphics[keepaspectratio]{bookdown-demo_files/figure-latex/unnamed-chunk-14-1.pdf}}

\begin{Shaded}
\begin{Highlighting}[]
\NormalTok{grafico\_mkt }\OtherTok{\textless{}{-}}\NormalTok{ grafico\_mkt }\SpecialCharTok{+} \FunctionTok{labs}\NormalTok{(}\AttributeTok{x =} \StringTok{"Estado civil"}\NormalTok{, }\AttributeTok{y =} \StringTok{"Recuento"}\NormalTok{)}
\NormalTok{grafico\_mkt}
\end{Highlighting}
\end{Shaded}

\pandocbounded{\includegraphics[keepaspectratio]{bookdown-demo_files/figure-latex/unnamed-chunk-14-2.pdf}}

\chapter{Representación de histogramas}\label{representacion_datos_histograma}

Podemos representar también los PIBs de \texttt{datos\_econ} y las edades de \texttt{datos\_mkt} en un histograma. Para ello, seguimos el mismo principio de antes pero solo necesitamos una variable (el PIB o la edad):

\begin{Shaded}
\begin{Highlighting}[]
\NormalTok{histograma\_econ }\OtherTok{\textless{}{-}} \FunctionTok{ggplot}\NormalTok{(datos\_econ, }\FunctionTok{aes}\NormalTok{(}\AttributeTok{x =}\NormalTok{ pib)) }\SpecialCharTok{+} \FunctionTok{geom\_histogram}\NormalTok{()}
\NormalTok{histograma\_econ}
\end{Highlighting}
\end{Shaded}

\begin{verbatim}
## `stat_bin()` using `bins = 30`. Pick better value `binwidth`.
\end{verbatim}

\pandocbounded{\includegraphics[keepaspectratio]{bookdown-demo_files/figure-latex/unnamed-chunk-15-1.pdf}}

\begin{Shaded}
\begin{Highlighting}[]
\NormalTok{histograma\_mkt }\OtherTok{\textless{}{-}} \FunctionTok{ggplot}\NormalTok{(datos\_mkt, }\FunctionTok{aes}\NormalTok{(}\AttributeTok{x =}\NormalTok{ edad)) }\SpecialCharTok{+} \FunctionTok{geom\_histogram}\NormalTok{()}
\NormalTok{histograma\_mkt}
\end{Highlighting}
\end{Shaded}

\begin{verbatim}
## `stat_bin()` using `bins = 30`. Pick better value `binwidth`.
\end{verbatim}

\pandocbounded{\includegraphics[keepaspectratio]{bookdown-demo_files/figure-latex/unnamed-chunk-15-2.pdf}}
La anchura de las barras que forman el histograma pueden ajustarse, y los títulos de los ejes se pueden cambiar de la misma forma que en el ejemplo anterior. Además, podemos darle color.

\begin{Shaded}
\begin{Highlighting}[]
\NormalTok{histograma\_econ }\OtherTok{\textless{}{-}} \FunctionTok{ggplot}\NormalTok{(datos\_econ, }\FunctionTok{aes}\NormalTok{(}\AttributeTok{x =}\NormalTok{ pib)) }\SpecialCharTok{+} \FunctionTok{geom\_histogram}\NormalTok{(}\AttributeTok{binwidth =} \DecValTok{300000}\NormalTok{, }\AttributeTok{fill =} \StringTok{"blue"}\NormalTok{)}
\NormalTok{histograma\_econ}
\end{Highlighting}
\end{Shaded}

\pandocbounded{\includegraphics[keepaspectratio]{bookdown-demo_files/figure-latex/unnamed-chunk-16-1.pdf}}

\begin{Shaded}
\begin{Highlighting}[]
\NormalTok{histograma\_mkt }\OtherTok{\textless{}{-}} \FunctionTok{ggplot}\NormalTok{(datos\_mkt, }\FunctionTok{aes}\NormalTok{(}\AttributeTok{x =}\NormalTok{ edad)) }\SpecialCharTok{+} \FunctionTok{geom\_histogram}\NormalTok{(}\AttributeTok{binwidth =} \DecValTok{2}\NormalTok{, }\AttributeTok{fill =} \StringTok{"red"}\NormalTok{)}
\NormalTok{histograma\_mkt}
\end{Highlighting}
\end{Shaded}

\pandocbounded{\includegraphics[keepaspectratio]{bookdown-demo_files/figure-latex/unnamed-chunk-16-2.pdf}}

Si queremos una paleta de colores más profesional, podemos usar el paquete \texttt{RColorBrewer}:

\begin{Shaded}
\begin{Highlighting}[]
\CommentTok{\# install.packages("RColorBrewer")}
\FunctionTok{library}\NormalTok{(RColorBrewer)}
\CommentTok{\# display.brewer.all() \# Muestra los colores disponibles}
\end{Highlighting}
\end{Shaded}

Para aplicar \texttt{RColorBrewer} podemos:

\begin{Shaded}
\begin{Highlighting}[]
\NormalTok{histograma\_econ }\OtherTok{\textless{}{-}} \FunctionTok{ggplot}\NormalTok{(datos\_econ, }\FunctionTok{aes}\NormalTok{(}\AttributeTok{x =}\NormalTok{ pib)) }\SpecialCharTok{+} 
  \FunctionTok{geom\_histogram}\NormalTok{(}\AttributeTok{binwidth =} \DecValTok{300000}\NormalTok{, }\AttributeTok{fill =} \FunctionTok{brewer.pal}\NormalTok{(}\DecValTok{5}\NormalTok{, }\StringTok{"Blues"}\NormalTok{)[}\DecValTok{3}\NormalTok{])}
\CommentTok{\# brewer.pal(n, name): n es el número de colores en la paleta y name es el nombre de la paleta, que puede consultarse en https://r{-}graph{-}gallery.com/38{-}rcolorbrewers{-}palettes.html. El [] posterior extrae el tercer color de la paleta. }
\NormalTok{histograma\_econ}
\end{Highlighting}
\end{Shaded}

\pandocbounded{\includegraphics[keepaspectratio]{bookdown-demo_files/figure-latex/unnamed-chunk-18-1.pdf}}

\begin{Shaded}
\begin{Highlighting}[]
\NormalTok{histograma\_mkt }\OtherTok{\textless{}{-}} \FunctionTok{ggplot}\NormalTok{(datos\_mkt, }\FunctionTok{aes}\NormalTok{(}\AttributeTok{x =}\NormalTok{ edad)) }\SpecialCharTok{+} 
  \FunctionTok{geom\_histogram}\NormalTok{(}\AttributeTok{binwidth =} \DecValTok{2}\NormalTok{, }\AttributeTok{fill =} \FunctionTok{brewer.pal}\NormalTok{(}\DecValTok{5}\NormalTok{, }\StringTok{"Reds"}\NormalTok{)[}\DecValTok{2}\NormalTok{])}
\CommentTok{\# brewer.pal(n, name): n es el número de colores en la paleta y name es el nombre de la paleta, que puede consultarse en https://r{-}graph{-}gallery.com/38{-}rcolorbrewers{-}palettes.html. El [] posterior extrae el tercer color de la paleta. }
\NormalTok{histograma\_mkt}
\end{Highlighting}
\end{Shaded}

\pandocbounded{\includegraphics[keepaspectratio]{bookdown-demo_files/figure-latex/unnamed-chunk-18-2.pdf}}
Podemos incluso superponer dos histogramas para observar si la distribución de la edad difiere entre el estado civil en \texttt{datos\_mkt}.

\begin{Shaded}
\begin{Highlighting}[]
\FunctionTok{ggplot}\NormalTok{(datos\_mkt, }\FunctionTok{aes}\NormalTok{(}\AttributeTok{x =}\NormalTok{ edad, }\AttributeTok{fill =}\NormalTok{ civil)) }\SpecialCharTok{+}
  \FunctionTok{geom\_histogram}\NormalTok{(}\AttributeTok{alpha =} \FloatTok{0.5}\NormalTok{, }\AttributeTok{position =} \StringTok{"identity"}\NormalTok{, }\AttributeTok{binwidth =} \DecValTok{3}\NormalTok{)}
\end{Highlighting}
\end{Shaded}

\pandocbounded{\includegraphics[keepaspectratio]{bookdown-demo_files/figure-latex/unnamed-chunk-19-1.pdf}}

\chapter{Representación de gráficos de líneas}\label{representaciuxf3n-de-gruxe1ficos-de-luxedneas}

Primero vamos a importar las base de datos de series temporales:

\begin{Shaded}
\begin{Highlighting}[]
\NormalTok{st\_econ }\OtherTok{\textless{}{-}} \FunctionTok{read\_excel}\NormalTok{(}\StringTok{"datos\_econ.xlsx"}\NormalTok{, }\DecValTok{3}\NormalTok{)}
\NormalTok{st\_econ }\CommentTok{\# Vemos qué contiene la base de datos}
\end{Highlighting}
\end{Shaded}

\begin{verbatim}
## # A tibble: 10 x 2
##     Year `Spain (GDP, M€)`
##    <dbl>             <dbl>
##  1  2016           1122967
##  2  2017           1170024
##  3  2018           1212276
##  4  2019           1253710
##  5  2020           1129214
##  6  2021           1235474
##  7  2022           1375863
##  8  2023           1497761
##  9  2024           1594330
## 10  2025           1685783
\end{verbatim}

\begin{Shaded}
\begin{Highlighting}[]
\FunctionTok{names}\NormalTok{(st\_econ) }\OtherTok{\textless{}{-}} \FunctionTok{c}\NormalTok{(}\StringTok{"año"}\NormalTok{, }\StringTok{"PIB"}\NormalTok{) }\CommentTok{\# cambiamos los nombres de las variables}
\NormalTok{st\_mkt }\OtherTok{\textless{}{-}}\NormalTok{ st\_mkt }\OtherTok{\textless{}{-}} \FunctionTok{read\_excel}\NormalTok{(}\StringTok{"datos\_mkt.xlsx"}\NormalTok{, }\DecValTok{2}\NormalTok{)}
\end{Highlighting}
\end{Shaded}

\begin{verbatim}
## New names:
## * `` -> `...1`
\end{verbatim}

\begin{Shaded}
\begin{Highlighting}[]
\NormalTok{st\_mkt }\CommentTok{\# Vemos qué contiene la base de datos}
\end{Highlighting}
\end{Shaded}

\begin{verbatim}
## # A tibble: 200 x 5
##     ...1    TV Radio Newspaper Sales
##    <dbl> <dbl> <dbl>     <dbl> <dbl>
##  1     1 230.   37.8      69.2  22.1
##  2     2  44.5  39.3      45.1  10.4
##  3     3  17.2  45.9      69.3   9.3
##  4     4 152.   41.3      58.5  18.5
##  5     5 181.   10.8      58.4  12.9
##  6     6   8.7  48.9      75     7.2
##  7     7  57.5  32.8      23.5  11.8
##  8     8 120.   19.6      11.6  13.2
##  9     9   8.6   2.1       1     4.8
## 10    10 200.    2.6      21.2  10.6
## # i 190 more rows
\end{verbatim}

\begin{Shaded}
\begin{Highlighting}[]
\FunctionTok{names}\NormalTok{(st\_mkt) }\OtherTok{\textless{}{-}} \FunctionTok{c}\NormalTok{(}\StringTok{"tiempo"}\NormalTok{, }\StringTok{"tv"}\NormalTok{, }\StringTok{"radio"}\NormalTok{, }\StringTok{"periodico"}\NormalTok{, }\StringTok{"ventas"}\NormalTok{) }\CommentTok{\# cambiamos los nombres de las variables}
\end{Highlighting}
\end{Shaded}

Ahora vamos a crear el gráfico usando los principios aprendidos anteriormente.

\begin{Shaded}
\begin{Highlighting}[]
\NormalTok{linea\_econ }\OtherTok{\textless{}{-}} \FunctionTok{ggplot}\NormalTok{(st\_econ, }\FunctionTok{aes}\NormalTok{(}\AttributeTok{x =}\NormalTok{ año, }\AttributeTok{y =}\NormalTok{ PIB)) }\SpecialCharTok{+}
  \FunctionTok{geom\_line}\NormalTok{(}\AttributeTok{color =} \FunctionTok{brewer.pal}\NormalTok{(}\DecValTok{3}\NormalTok{, }\StringTok{"Reds"}\NormalTok{)[}\DecValTok{3}\NormalTok{]) }\SpecialCharTok{+}
  \FunctionTok{labs}\NormalTok{(}\AttributeTok{y =} \StringTok{"Producto interior bruto (en millones de €) de España"}\NormalTok{, }\AttributeTok{x =} \StringTok{"Año"}\NormalTok{)}
\NormalTok{linea\_econ}
\end{Highlighting}
\end{Shaded}

\pandocbounded{\includegraphics[keepaspectratio]{bookdown-demo_files/figure-latex/unnamed-chunk-21-1.pdf}}

\begin{Shaded}
\begin{Highlighting}[]
\NormalTok{linea\_mkt }\OtherTok{\textless{}{-}} \FunctionTok{ggplot}\NormalTok{(st\_mkt, }\FunctionTok{aes}\NormalTok{(}\AttributeTok{x =}\NormalTok{ tiempo, }\AttributeTok{y =}\NormalTok{ ventas)) }\SpecialCharTok{+}
  \FunctionTok{geom\_line}\NormalTok{(}\AttributeTok{color =} \FunctionTok{brewer.pal}\NormalTok{(}\DecValTok{3}\NormalTok{, }\StringTok{"Reds"}\NormalTok{)[}\DecValTok{3}\NormalTok{]) }\SpecialCharTok{+}
  \FunctionTok{labs}\NormalTok{(}\AttributeTok{x =} \StringTok{"Tiempo"}\NormalTok{, }\AttributeTok{y =} \StringTok{"Ventas"}\NormalTok{)}
\NormalTok{linea\_mkt}
\end{Highlighting}
\end{Shaded}

\pandocbounded{\includegraphics[keepaspectratio]{bookdown-demo_files/figure-latex/unnamed-chunk-21-2.pdf}}

Incluso en vez de \texttt{geom\_line}, puede usarse \texttt{geom\_smooth} para suavizar la línea:

\begin{Shaded}
\begin{Highlighting}[]
\NormalTok{linea\_econ }\OtherTok{\textless{}{-}} \FunctionTok{ggplot}\NormalTok{(st\_econ, }\FunctionTok{aes}\NormalTok{(}\AttributeTok{x =}\NormalTok{ año, }\AttributeTok{y =}\NormalTok{ PIB)) }\SpecialCharTok{+}
  \FunctionTok{geom\_smooth}\NormalTok{(}\AttributeTok{color =} \FunctionTok{brewer.pal}\NormalTok{(}\DecValTok{3}\NormalTok{, }\StringTok{"Blues"}\NormalTok{)[}\DecValTok{3}\NormalTok{]) }\SpecialCharTok{+}
  \FunctionTok{labs}\NormalTok{(}\AttributeTok{y =} \StringTok{"Producto interior bruto (en millones de €) de España"}\NormalTok{, }\AttributeTok{x =} \StringTok{"Año"}\NormalTok{)}
\NormalTok{linea\_econ}
\end{Highlighting}
\end{Shaded}

\begin{verbatim}
## `geom_smooth()` using method = 'loess' and formula = 'y ~ x'
\end{verbatim}

\pandocbounded{\includegraphics[keepaspectratio]{bookdown-demo_files/figure-latex/unnamed-chunk-22-1.pdf}}

\begin{Shaded}
\begin{Highlighting}[]
\NormalTok{linea\_mkt }\OtherTok{\textless{}{-}} \FunctionTok{ggplot}\NormalTok{(st\_mkt, }\FunctionTok{aes}\NormalTok{(}\AttributeTok{x =}\NormalTok{ tiempo, }\AttributeTok{y =}\NormalTok{ ventas)) }\SpecialCharTok{+}
  \FunctionTok{geom\_smooth}\NormalTok{(}\AttributeTok{color =} \FunctionTok{brewer.pal}\NormalTok{(}\DecValTok{3}\NormalTok{, }\StringTok{"Reds"}\NormalTok{)[}\DecValTok{3}\NormalTok{]) }\SpecialCharTok{+}
  \FunctionTok{labs}\NormalTok{(}\AttributeTok{x =} \StringTok{"Tiempo"}\NormalTok{, }\AttributeTok{y =} \StringTok{"Ventas"}\NormalTok{)}
\NormalTok{linea\_mkt}
\end{Highlighting}
\end{Shaded}

\begin{verbatim}
## `geom_smooth()` using method = 'loess' and formula = 'y ~ x'
\end{verbatim}

\pandocbounded{\includegraphics[keepaspectratio]{bookdown-demo_files/figure-latex/unnamed-chunk-22-2.pdf}}

\bibliography{book.bib,packages.bib}

\end{document}
